% ============================================================
% 11_math_advanced.tex — OPTIONALES Math-Erweiterungsmodul
% ============================================================
% Opt-in: in preamble.tex einkommentieren für mathe-lastige
% Fächer (LinAlg, Analysis). Baut auf 10_math.tex auf und
% ergänzt nur das Delta — keine Basis-Makros duplizieren.
%
% Voraussetzungen (aus 00_packages.tex):
%   - mathtools, unicode-math   (Operatoren, Symbole)
%   - tikz                      (drawbrace / tikzmark / annote)
% ============================================================

% --- Lineare-Algebra-Operatoren ---
\DeclareMathOperator{\Ker}{Ker}
\DeclareMathOperator{\rang}{rang}
\DeclareMathOperator{\Spur}{Spur}
\DeclareMathOperator{\tr}{tr}
\DeclareMathOperator{\diag}{diag}
\DeclareMathOperator{\spanop}{span}
\DeclareMathOperator{\eig}{eig}
\DeclareMathOperator{\algVfh}{algVfh}
\DeclareMathOperator{\gVfh}{gVfh}
\DeclareMathOperator{\proj}{proj}
\DeclareMathOperator{\sig}{sig}
\DeclareMathOperator{\Adj}{Adj}

% --- Analysis-Operatoren (Abweichung dieses Forks) ---
% Grad(p) ist der Polynomgrad und bewusst gross geschrieben — \grad aus
% 10_math ist der Gradient und eine andere Sache.
% Die Areafunktionen sind die Umkehrungen der Hyperbelfunktionen; TeX bringt
% dafuer nichts mit, und \operatorname im Kapitel meldet der Linter zu Recht:
% Ein Operator, der mehrfach vorkommt, gehoert einmal hierher.
\DeclareMathOperator{\Grad}{Grad}
\DeclareMathOperator{\Arsinh}{Arsinh}
\DeclareMathOperator{\Arcosh}{Arcosh}
\DeclareMathOperator{\Artanh}{Artanh}

% --- Phasenportraet / GGWP Semantik-Typen (Ink-konform) ---
\newcommand{\ZSFggwpKnoten}[1][Knoten]{\ZSFInk{ChapterColor9}{\textbf{#1}}}
\newcommand{\ZSFggwpSattel}[1][Sattel]{\ZSFInk{ChapterColor14}{\textbf{#1}}}
\newcommand{\ZSFggwpStrudel}[1][Strudel]{\ZSFInk{ChapterColor4}{\textbf{#1}}}
\newcommand{\ZSFggwpZentrum}[1][Zentrum]{\textbf{#1}}

% --- Gotisches Im/Re -> moderne aufrechte Operatoren ---
% AtBeginDocument, weil unicode-math seine Symbole ebenfalls erst dort
% definiert — eine Praeambel-Renewcommand wuerde still ueberschrieben.
\AtBeginDocument{%
  \let\oldIm\Im
  \renewcommand{\Im}{\operatorname{Im}}%
  \let\oldRe\Re
  \renewcommand{\Re}{\operatorname{Re}}%
}

% --- TikZ-Klammern für Matrix-Annotationen ---
\usetikzlibrary{decorations.pathreplacing,calc}
% Das optionale Argument sind KNOTEN-OPTIONEN (z.B. xshift), keine Länge. Der
% Default ist deshalb leer: Ein Mass an dieser Stelle liest pgf als
% Pfeilspezifikation und bricht ab ("Unknown arrow tip kind"). Der vertikale
% Versatz der Klammer hängt ohnehin an den Koordinaten, siehe unten.
\newcommand{\tikzmark}[2][]{\tikz[remember picture, overlay, baseline=-0.5ex]\node[#1](#2){};}
% Die Marker von \tikzmark sitzen auf der GRUNDLINIE des Terms. Ohne Versatz
% liegt die Klammer-Sehne genau dort und läuft als waagrechte Linie mitten
% durch die Formel. Der Versatz muss an den Koordinaten hängen, nicht als
% yshift an der Pfadoption: Knoten aus einem anderen Bild (remember picture)
% werden absolut aufgelöst und ignorieren die Pfad-Transformation.
%
% Richtung und Dekoration gehören zusammen, deshalb setzt jeder Stil beides.
% \ZSF@braceSign trägt das Vorzeichen zur Koordinate — die Optionen werden vor
% den Koordinaten ausgewertet, der Wert steht also rechtzeitig.
\makeatletter
\newcommand{\ZSF@braceSign}{-}
\tikzset{
  brace/.style={decorate, decoration={brace},
                /utils/exec=\renewcommand{\ZSF@braceSign}{}},
  brace mirrored/.style={decorate, decoration={brace,mirror},
                /utils/exec=\renewcommand{\ZSF@braceSign}{-}},
}
\newcounter{brace}
\setcounter{brace}{0}
% \drawbrace[stil]{marke links}{marke rechts}
%   brace           — Klammer ÜBER dem Term (Default)
%   brace mirrored  — Klammer UNTER dem Term
\newcommand{\drawbrace}[3][brace]{%
  \refstepcounter{brace}%
  \tikz[remember picture, overlay]\draw[#1]
    ([yshift=\ZSF@braceSign\ZSFbraceDrop]#2.center)--%
    ([yshift=\ZSF@braceSign\ZSFbraceDrop]#3.center)%
    node[pos=0.5, name=brace-\thebrace]{};%
}
\makeatother
\newcommand{\annote}[3][]{%
  \tikz[remember picture, overlay]\node[#1] at (#2) {#3};%
}

%
%
%
%
%
%
%
%
%
%
\newcommand{\ZSFbraceRoom}{\vspace{\dimexpr\ZSFbraceDrop+\ZSFspaceM\relax}}

% --- Asterisk-Definitionen (mathabx-Namen auf Standard-Symbole) ---
\newcommand{\asterisk}{\ast}
\newcommand{\oasterisk}{\circledast}
